\begin{table}[htbp]
\centering
\normalsize
\setlength{\tabcolsep}{4pt}
\renewcommand{\arraystretch}{1.12}
\begin{tabularx}{\textwidth}{L{4.0cm} Y}
\hline
\textbf{Configuration item} & \textbf{Implementation report} \\
\hline
drag\&bot platform &
drag\&bot runtime package 3.6.0, used as the virtual robot-programming and
execution environment for the thesis implementation. \\
\hline
Runtime interface &
AnyRobot WebSocket-based interface used for command exchange between the
planning software and the drag\&bot runtime. \\
\hline
Simulation setup &
Virtual press-brake cell containing the UR5e robot model, press brake,
orientation table, flip station, vacuum gripper, and tracked sheet object. \\
\hline
Robot configuration &
UR5e manipulator model with six degrees of freedom, used for
all simulated handling and bending-related motion sequences. \\
\hline
\ac{tcp} definition &
Configured TCP \texttt{project\_3\_tcp} with offsets
$x = 0{,}000$~m, $y = 0{,}034$~m, $z = 0{,}194$~m,
$RZ = -90{,}0^\circ$, $RY = -75{,}0^\circ$, and $RX = 0{,}0^\circ$. \\
\hline
Gripper setup &
Two-cup vacuum gripper represented in the simulation by digital input-output
control, vacuum activation and release states, and kinematic object
attachment-detachment logic. \\
\hline
System configuration &
Workpiece dimensions, bend-side selection, safety limits, admissible workspace
zones, gripper parameters, and runtime settings are loaded from configuration
data before execution. \\
\hline
\end{tabularx}
\caption{Hardware and platform configuration used for the simulation-based implementation.}
\label{tab:hardware-implementation-report}
\end{table}
